\section{FPS Analysis}
In this section we present our \textit{Fixed-Priority Scheduling (FPS)} analysis
of the system. The analysis was performed by modeling the system through a
\textit{MAST} model, which was then analyzed using the \textit{MAST} tool to verify
the timing properties and to analyze the system under different workloads. \\
The priority assignment was done according to the \textit{Deadline Monotonic} policy, which
assigns higher priorities to tasks with shorter deadlines. This choice comes
from the fact that tasks in the system do not have implicit deadlines and, as
theory proves, \textit{Deadline Monotonic} performs better than \textit{Rate  
Monotonic} in such cases, as it can provide feasible schedules where
\textit{Rate Monotonic} fails. \\
The evaluation of the system under different utilization levels was facilitated
by the use of the \textit{Small Whetstone} algorithm to simulate computational workloads.
Thanks to its configurable load and deterministic behaviour without involving
the use of shared resources, it provides a straightforward way to test the
system under differen utilizations by simply scaling the WCET in the different
\texttt{small\_whetstone} operations that are present in the \textit{MAST}
model, without the need for new execution-time measurements.  

\subsection{Holistic Analysis}
We first conducted an holistic analysis, considering all the tasks as
intendepent. As a consequence, this analysis is pessimistic, since it does not
take into account the relationships between tasks. System feasibility is
therefore evaluated based on the worst-case scenario at the critical instant, i.e.,
when all tasks are released simultaneously (not a realistic and possible case).
The analysis was performed using the \textit{Classic RM Analysis} tool in
\textit{MAST}, which implements the classic exact response-time analysis for
single-processor systems under \textit{Fixed-Priority Scheduling} (although it's
called \textit{Rate Monotonic}, it bases \textit{Deadline Monotonic} analysis). \\
The first results of the analysis were made considering the system under the
loads provided in the original \textit{Guided Extended Example}, i.e.,
with the following Whetstone workloads (accompained by their corresponding WCETs
expressed in seconds).
\begin{table}[H]
  \label{tab:holistic-analysis-ada}
  \resizebox{\linewidth}{!}{%
    \begin{tabular}{lccc}
      \toprule
      \textbf{Task} & \textbf{Workload} & \textbf{Rust WCET} & \textbf{Ada WCET} \\ 
      \midrule
      \texttt{Regular\_Producer}       & 756  & 0.023639364 & 0.00554170 \\
      \texttt{On\_Call\_Producer}      & 278  & 0.008692782 & 0.00203782 \\
      \texttt{Activation\_Log\_Reader} & 139  & 0.004346391 & 0.00101891 \\
      \bottomrule
    \end{tabular}
  }
  \caption{Whetstone Operations}
\end{table}
The obtained results are reported in Table [\ref{tab:holistic-analysis-rust}] and
Table [\ref{tab:holistic-analysis-ada}]. Times are expressed in seconds and we
omitted the watchdogs' transactions (which were instead included in the Rust
model) since they are not relevant for identifying increasing workload factors.
Regarding the jitter times, all transactions suffer jitter due to the system
ticker interrupt and the context switch overhead. Moreover, additional sources 
of jitter and blocking times for each transaction are the following:
\begin{itemize}
  \item \texttt{rp\_transaction}: suffers additional jitter due to the possible
  interference of the \texttt{ees\_transaction}. Its blocking time is instead due
  to the possible access to shared resources from the lower-priority
  \texttt{ocp\_transaction} and \texttt{alr\_transaction} with ceiling
  priority equal to the priority of \texttt{rp\_transaction};
  \item \texttt{ocp\_transaction}: suffers additional jitter due to interference
  by the higher-priority \texttt{rp\_transaction}. Its blocking time is caused
  by the \texttt{alr\_transaction};
  \item \texttt{alr\_transaction}: suffers additional jitter due to interference
  by both the higher-priority \texttt{rp\_transaction} and
  \texttt{ocp\_transaction}. It does not suffer blocking, since it represents
  the lowest-priority task in the system;
  \item \texttt{ees\_transaction}: it suffers no additional jitter than the
  aforementioned overheads. Its blockig time is caused by the
  \texttt{alr\_transaction} which can possibly access the shared
  \texttt{Activation\_Log} resource with ceiling priority equal to the
  priority of \texttt{ees\_transaction}.
\end{itemize}
\begin{table}[H]
  \resizebox{\linewidth}{!}{%
    \begin{tabular}{lcccc}
      \toprule
      \textbf{Transaction} & \boldmath{$R_{max}$} & \textbf{Slack} &
      \textbf{Jitter} & \textbf{MBT} \\ 
      \midrule
      \texttt{rp\_transaction}  & 0.024728  & 2006.3\%     & 0.001066  & 3.030E-06 \\
      \texttt{ocp\_transaction} & 0.033451  & 8791.0\%     & 0.024741  & 2.730E-06 \\
      \texttt{alr\_transaction} & 0.037822  & 21927.3\%    & 0.033455  & 0.000     \\
      \texttt{ees\_transaction} & 0.001043  & >=100000.0\% & 2.955E-05 & 2.730E-06 \\
      \hline
      \texttt{System Utilization} & 2.79\% & & & \\
      \texttt{System Slack} & 1994.1\% & & & \\
      \bottomrule
    \end{tabular}
  }
  \caption{Holistic Analysis Results Rust}
  \label{tab:holistic-analysis-rust}
\end{table}
Based on these results, since the smallest slack available was 2006.3\% for the 
\texttt{rp\_transaction}, we increased the workloads with a factor of 20 to
reduce it significantly. New results are reported in Table
[\ref{tab:holistic-analysis-rust-x20}]. 
\begin{table}[H]
  \resizebox{\linewidth}{!}{%
    \begin{tabular}{lcccc}
      \toprule
      \textbf{Transaction} & \boldmath{$R_{max}$} & \textbf{Slack} &
      \textbf{Jitter} & \textbf{MBT} \\ 
      \midrule
      \texttt{rp\_transaction}  & 0.474326  & 5.08\%       & 0.001516  & 3.030E-06 \\
      \texttt{ocp\_transaction} & 0.648377  & 87.11\%      & 0.474643  & 2.730E-06 \\
      \texttt{alr\_transaction} & 0.735413  & 303.91\%     & 0.787757  & 0.000     \\
      \texttt{ees\_transaction} & 0.001043  & >=100000.0\% & 0.001030  & 2.730E-06 \\
      \hline
      \texttt{System Utilization} & 53.76\% & & & \\
      \texttt{System Slack} & 5.75\% & & & \\
      \bottomrule
    \end{tabular}
  }
  \caption{Holistic Analysis Results Rust x20}
  \label{tab:holistic-analysis-rust-x20}
\end{table}
We can notice how the blocking times have not changed, since protected operations
remain the same. While we have achieved a small slack for the
\texttt{rp\_transaction} (5.08\%), the other two transactions that involve
Whetstone computations still have large available slacks. Based on these
results, especially on the smallest slack of 87.11\%, we decided to increase the
\texttt{On\_Call\_Producer} and \\
\texttt{Activation\_Log\_Reader} workloads by a factor 1.8, while keeping the
\texttt{Regular\_Producer} workload fixed. The updated results are reported in
Table [\ref{tab:holistic-analysis-rust-x1.8}].  
\begin{table}[H]
  \resizebox{\linewidth}{!}{%
    \begin{tabular}{lcccc}
      \toprule
      \textbf{Transaction} & \boldmath{$R_{max}$} & \textbf{Slack} &
      \textbf{Jitter} & \textbf{MBT} \\ 
      \midrule
      \texttt{rp\_transaction}  & 0.474326  & 2.34\%      & 0.001516  & 3.030E-06 \\
      \texttt{ocp\_transaction} & 0.787600  & 3.91\%      & 0.473643  & 2.730E-06 \\
      \texttt{alr\_transaction} & 0.944248  & 35.55\%     & 0.786757  & 0.000     \\
      \texttt{ees\_transaction} & 0.001043  & 91855.5\%   & 0.001030  & 2.730E-06 \\
      \hline
      \texttt{System Utilization} & 58.86\% & & & \\
      \texttt{System Slack} & 1.58\% & & & \\
      \bottomrule
    \end{tabular}
  }
  \caption{Holistic Analysis Results Rust x1.8}
  \label{tab:holistic-analysis-rust-x1.8}
\end{table}
Finally, based on the last useful available slack of 35.55\% for the
\texttt{alr\_transaction}, we increased its workload by a factor of 1.35. The
updated results are reported in Table [\ref{tab:holistic-analysis-rust-x1.35}].
\begin{table}[H]
  \resizebox{\linewidth}{!}{%
    \begin{tabular}{lcccc}
      \toprule
      \textbf{Transaction} & \boldmath{$R_{max}$} & \textbf{Slack} &
      \textbf{Jitter} & \textbf{MBT} \\ 
      \midrule
      \texttt{rp\_transaction}  & 0.474326  & 0.00\%      & 0.001516  & 3.030E-06 \\
      \texttt{ocp\_transaction} & 0.787600  & 0.00\%      & 0.474643  & 2.730E-06 \\
      \texttt{alr\_transaction} & 0.999068  & 0.39\%      & 0.787812  & 0.000     \\
      \texttt{ees\_transaction} & 0.001043  & 6911.7\%    & 0.001030  & 2.730E-06 \\
      \hline
      \texttt{System Utilization} & 60.68\% & & & \\
      \texttt{System Slack} & 0.39\% & & & \\
      \bottomrule
    \end{tabular}
  }
  \caption{Holistic Analysis Results Rust x1.35}
  \label{tab:holistic-analysis-rust-x1.35}
\end{table}
Even though the \texttt{ees\_transaction} has still a large available slack,
since it does not perform any Whetstone computation, it is not possible to 
increase the utilization of the system any further. Besides that, even an
eventual increase of the execution time of \texttt{ees\_transaction} would
not affect significantly the utilization of the system. In conclusion, under an
holistic analysis, we achieved an utilization of the system equal to 60.68\%.

\begin{table}[H]
  \resizebox{\linewidth}{!}{%
    \begin{tabular}{lcccc}
      \toprule
      \textbf{Transaction} & \boldmath{$R_{max}$} & \textbf{Slack} &
      \textbf{Jitter} & \textbf{MBT} \\ 
      \midrule
      \texttt{rp\_transaction}  & 0.006624  & 8808.2\%     & 0.001045  & 1.361E-06 \\
      \texttt{ocp\_transaction} & 0.007697  & 38006.3\%    & 0.005638  & 1.267E-06 \\
      \texttt{alr\_transaction} & 0.008752  & 91394.9\%    & 0.007706  & 0.000     \\
      \texttt{ees\_transaction} & 1.882E-05 & >=100000.0\% & 1.141E-05 & 1.267E-06 \\
      \hline
      \texttt{System Utilization} & 1.03\% & & & \\
      \texttt{System Slack} & 7728.7\% & & & \\
      \bottomrule
    \end{tabular}
  }
  \caption{Holistic Analysis Results Ada}
  \label{tab:holistic-analysis-ada}
\end{table}
Based on these results, since the smallest available slack was 8808.2\% for the 
\texttt{rp\_transaction}, we increased the workloads with a factor of 90 to
reduce it significantly. The updated results are reported in Table
[\ref{tab:holistic-analysis-ada-x88}].  
\begin{table}[H]
  \resizebox{\linewidth}{!}{%
    \begin{tabular}{lcccc}
      \toprule
      \textbf{Transaction} & \boldmath{$R_{max}$} & \textbf{Slack} &
      \textbf{Jitter} & \textbf{MBT} \\ 
      \midrule
      \texttt{rp\_transaction}  & 0.490670  & 1.56\%       & 0.002963  & 1.361E-06 \\
      \texttt{ocp\_transaction} & 0.669740  & 72.27\%      & 0.490390  & 1.267E-06 \\
      \texttt{alr\_transaction} & 0.759792  & 265.23\%     & 0.670102  & 0.000     \\
      \texttt{ees\_transaction} & 1.882E-05 & >=100000.0\% & 1.141E-05 & 1.267E-06 \\
      \hline
      \texttt{System Utilization} & 55.74\% & & & \\
      \texttt{System Slack} & 1.98\% & & & \\
      \bottomrule
    \end{tabular}
  }
  \caption{Holistic Analysis Results Ada x88}
  \label{tab:holistic-analysis-ada-x88}
\end{table}
Again, we can notice how the blocking times have not changed, since protected
operations remain the same. While we have achieved a small slack for the
\texttt{rp\_transaction} (1.56\%), the other two transactions that involve
Whetstone computations still have large available slacks. Based on these
results, especially on the smallest slack of 72.27\%, we decided to increase the
\texttt{On\_Call\_Producer} and \\
\texttt{Activation\_Log\_Reader} workloads by a factor of 1.7, while keeping the
\texttt{Regular\_Producer} workload fixed. The updated results are reported in Table
[\ref{tab:holistic-analysis-ada-x1.7}].  
\begin{table}[H]
  \resizebox{\linewidth}{!}{%
    \begin{tabular}{lcccc}
      \toprule
      \textbf{Transaction} & \boldmath{$R_{max}$} & \textbf{Slack} &
      \textbf{Jitter} & \textbf{MBT} \\ 
      \midrule
      \texttt{rp\_transaction}  & 0.490670  & 0.78\%       & 0.002963  & 1.361E-06 \\
      \texttt{ocp\_transaction} & 0.795769  & 1.17\%       & 0.490889  & 0.000     \\
      \texttt{alr\_transaction} & 0.948836  & 32.42\%      & 0.796380  & 1.361E-06 \\
      \texttt{ees\_transaction} & 1.882E-05 & 56921.9\%    & 1.141E-05 & 1.267E-06 \\
      \hline
      \texttt{System Utilization} & 60.35\% & & & \\
      \texttt{System Slack} & 0.78\% & & & \\
      \bottomrule
    \end{tabular}
  }
  \caption{Holistic Analysis Results Ada x1.7}
  \label{tab:holistic-analysis-ada-x1.7}
\end{table}
Finally, based on the last useful slack available of 32.34\% for the
\texttt{alr\_transaction}, we increased its workload by a factor of 1.3. The
updated results are reported in Table [\ref{tab:holistic-analysis-ada-x1.3}].
\begin{table}[H]
  \resizebox{\linewidth}{!}{%
    \begin{tabular}{lcccc}
      \toprule
      \textbf{Transaction} & \boldmath{$R_{max}$} & \textbf{Slack} &
      \textbf{Jitter} & \textbf{MBT} \\ 
      \midrule
      \texttt{rp\_transaction}  & 0.490670  & 0.78\%       & 0.002963  & 1.361E-06 \\
      \texttt{ocp\_transaction} & 0.795769  & 1.17\%       & 0.490889  & 0.000     \\
      \texttt{alr\_transaction} & 0.994747  & 1.95\%       & 0.796563  & 1.361E-06 \\
      \texttt{ees\_transaction} & 1.882E-05 & 56921.9\%    & 1.141E-05 & 1.267E-06 \\
      \hline
      \texttt{System Utilization} & 61.87\% & & & \\
      \texttt{System Slack} & 0.78\% & & & \\
      \bottomrule
    \end{tabular}
  }
  \caption{Holistic Analysis Results Ada x1.3}
  \label{tab:holistic-analysis-ada-x1.3}
\end{table}
Again, the considerations made before upon the \texttt{ees\_transaction} apply
here. In conclusion, under an holistic analysis, we achieved an utilization of the system
equal to 61.87\%. 


